\section*{Chapter \ref{ch:ml:training-infrastructure} --- Training Infrastructure}

\begin{solution}{exo:ml:training-infrastructure:1}
\omterm{def:ml:training-infrastructure:mixed}{bfloat16} $2^{-7}\approx7.8\times10^{-3}$, float16 $2^{-10}\approx9.8\times10^{-4}$, float32 $2^{-23}\approx1.2\times10^{-7}$.
$10^{-30}$ is representable in \omterm{def:ml:training-infrastructure:mixed}{bfloat16} and float32 (smallest normal numbers about $10^{-38}$) but not in float16, whose
smallest positive number is about $6\times10^{-8}$: it underflows to zero, which is why float16 training needs loss scaling.
\end{solution}

\begin{solution}{exo:ml:training-infrastructure:2}
$1\,140 - 20 - 10 + 1 = 1\,111$ windows. Uniform starts over the concatenation would produce windows (or targets) that begin in one
session and end in the next, splicing a close to an open: a price jump the model would learn to predict from nothing.
\end{solution}

\begin{solution}{exo:ml:training-infrastructure:3}
With the memory map and the loop, a step is $0.74 + 1.17 = 1.91$~ms; halving compute gives 1.32~ms, a speed-up of 1.44. With the
gather, $0.27 + 1.10 = 1.37$~ms becomes 0.82~ms, a speed-up of 1.67. The faster the loader, the more a faster device is worth.
\end{solution}

\begin{solution}{exo:ml:training-infrastructure:4}
Each step's rounding changes the update slightly; the two trajectories diverge in parameter space and see the same batches with
different parameters, so their per-step losses differ by as much as the loss varies from batch to batch. Both converge to
models of the same quality; the comparison must be of the metric over the run, not step by step.
\end{solution}

\begin{solution}{exo:ml:training-infrastructure:5}
The loss on the union is $\frac1{kb}\sum_{j=1}^k\sum_{i\in B_j}\ell_i = \frac1k\sum_jL_j$ with $L_j$ each micro-batch's mean, and
the gradient is linear: summing the gradients of $L_j/k$ gives it. \omterm{def:ml:neural-networks-for-noisy-tabular-data:dropout}{Batch normalisation} computes statistics per micro-batch, so the
forward pass itself differs and the equivalence breaks; so does anything else that couples examples within a batch.
\end{solution}

\begin{solution}{exo:ml:training-infrastructure:6}
Weights alone do not determine the rest of the run: the optimiser's moments restart, the sampler replays or skips data, and the
random generators restart. In the chapter's experiment the resumed run ended up to 0.039 away from the uninterrupted one. Save and
restore every state, and test resumption bitwise.
\end{solution}

\begin{solution}{exo:ml:training-infrastructure:7}
Rows: $250\times6.5\times3\,600/0.5 = 11.7$ million; binary $11.7\times10^6\times41\times4 = 1.92$~GB; text about 2.8~GB. At 176~MB/s
the text takes about 16~s per \omterm{def:ml:neural-networks-for-noisy-tabular-data:backprop}{epoch}; the NumPy files at 1.8~GB/s about 1.1~s; the memory map reads only what the batches touch,
from the page cache once warm.
\end{solution}

\begin{solution}{exo:ml:training-infrastructure:8}
In a ring, the array is cut into $p$ chunks; a reduce-scatter passes chunks round the ring $p - 1$ times, each device sending
$n/p$ numbers each time, and an all-gather does the same again: $2n(p-1)/p$ in all, nearly $2n$ whatever $p$. The communication
per device does not grow with the number of devices, but it does not shrink either, while each device's compute shrinks as $1/p$:
beyond some $p$, communication dominates.
\end{solution}

\begin{solution}{pb:ml:training-infrastructure:1}
\textbf{Part I.}
\begin{enumerate}
\item CSV 2.21~MB, NumPy 1.50~MB, memory map 1.50~MB.
\item Load per \omterm{def:ml:neural-networks-for-noisy-tabular-data:backprop}{epoch} 12.5~ms (CSV), 0.83~ms (NumPy), 0.39~ms (map); batch 0.48 to 0.74~ms looped and 0.27~ms gathered; compute
1.1 to 1.2~ms (measured on the author's laptop).
\item 0.34 (CSV), 0.31 (NumPy), 0.39 (map, looped) and 0.20 (map, gathered).
\item The batch assembly: a vectorised gather halves the data's share before any storage change matters.
\end{enumerate}
\textbf{Part II.}
\begin{enumerate}[resume]
\item Eight exponent bits, seven fraction bits; float32's range; epsilon $2^{-7}$.
\item Last-50-step mean losses 0.0525 (\omterm{def:ml:training-infrastructure:mixed}{bfloat16}) and 0.0527 (float32); largest step difference 0.0132 (22\%).
\item Its exponent range equals float32's: small gradients do not underflow.
\item Loss, optimiser state and weights; accumulations of many terms (Book~4's compensated sums); risk and P\&L computations
downstream of the model.
\end{enumerate}
\textbf{Part III.}
\begin{enumerate}[resume]
\item Accumulation equals a larger batch to $1.4\times10^{-7}$; averaged shard gradients equal the full gradient to
$1.5\times10^{-8}$.
\item Model, optimiser, sampler and generator states and the step; with weights alone the run ended up to 0.039 away.
\item 693~kB.
\item Rescaling with the batch, usually with a warm-up; and a check that the larger batch improves the model at all.
\end{enumerate}
\textbf{Part IV.}
\begin{enumerate}[resume]
\item \emph{Where the time goes.} The data take 0.34, 0.31 and 0.39 of a training step with CSV, NumPy and memory-mapped storage and
window-by-window batching, and 0.20 with the memory map and one gather (laptop measurements); the largest loss difference between
float32 and \omterm{def:ml:training-infrastructure:mixed}{bfloat16} over the 300-step run is 0.0132, while their final losses agree to 0.0002.
\item Not before the loader: a faster device would wait for data a third of the time.
\item When the model's parameters and activations do not fit on one device: rarely for trading models, often for \omterm{def:ml:large-language-models-in-finance:lm}{language models}.
\item Fix seeds, deterministic algorithms, library versions and the data order; accept that different hardware can round
differently and compare on metrics.
\item Windows spanning a close and an open teach the model a jump it cannot predict and leak information across days.
\item Many repetitions, medians, the machine's load recorded, and ratios rather than absolute times.
\item Code version, data version, configuration, seeds, the machine, the checkpoints and the metrics (chapter~25).
\item The share of each step spent waiting for data.
\end{enumerate}
\end{solution}

\begin{solution}{iq:ml:training-infrastructure:1}
A 16-bit format with float32's exponent and seven fraction bits: less precision than float16 but the full range, so gradients do
not underflow and no loss scaling is needed; matrix units run it at full speed.
\iqlookfor{the bit layout and the range argument.}
\end{solution}

\begin{solution}{iq:ml:training-infrastructure:2}
Profile the step: time data loading, host-to-device transfer and compute separately; check loader parallelism and prefetching,
batch size, synchronisation points (logging, .item() calls), and small kernels; fix the largest.
\iqlookfor{measurement first, then the usual culprits.}
\end{solution}

\begin{solution}{iq:ml:training-infrastructure:3}
Each device holds a copy of the model and processes its share of the batch; after the backward pass an \omterm{def:ml:training-infrastructure:dist}{all-reduce} leaves every
device with the average gradient, and all apply the same update.
\iqlookfor{replicas, sharded batches, averaged gradients.}
\end{solution}

\begin{solution}{iq:ml:training-infrastructure:4}
Per-instrument, per-session binary columns (memory-mapped or chunked), an index of session boundaries, samplers that draw windows
within sessions and respect the train/validation split in time, batch gathering in one read, and normalisation fitted on training
data only.
\iqlookfor{binary storage, session-aware sampling and time-respecting splits.}
\end{solution}

\begin{solution}{iq:ml:training-infrastructure:5}
Model parameters and buffers, optimiser state, learning-rate schedule state, data sampler position and random state, all random
generator states, the step or \omterm{def:ml:neural-networks-for-noisy-tabular-data:backprop}{epoch} count, and the code and data versions.
\iqlookfor{the optimiser and the sampler, not only the weights.}
\end{solution}

\begin{solution}{iq:ml:training-infrastructure:6}
When the model has per-batch statistics (\omterm{def:ml:neural-networks-for-noisy-tabular-data:dropout}{batch normalisation}), when the micro-batches are not equal in size, when the order of
floating-point summation matters for bitwise results, and when devices see different data distributions.
\iqlookfor{batch-coupled layers, unequal shards and summation order.}
\end{solution}
